my% ============================================================================
% ISBA 2026 POSTER — HIGH-AESTHETIC, SELF-CONTAINED BEAMERPOSTER SOURCE
% Compile with pdfLaTeX on Overleaf.
%
% CLAIM GOVERNANCE — DO NOT CHANGE WITHOUT PROMOTED/FROZEN RESULTS:
%   1. exAL-M-T1 has the lowest 28-day CRPS at 4 of 5 rolling origins.
%   2. AL-M-T1 has the lowest 28-day CRPS at the 2022-12-25 origin.
%   3. Screening/discount/epsilon runs are non-authoritative and must not appear.
% ============================================================================

\documentclass[final]{beamer}

\usepackage[
  size=custom,
  width=84.1,
  height=118.9,
  scale=1.0
]{beamerposter}

% ---------- Typography -------------------------------------------------------
% These packages are available in standard Overleaf TeX Live installations.
\usepackage[T1]{fontenc}
\usepackage[utf8]{inputenc}
\usepackage[type1,defaultsans]{lato}
\usepackage[type1,nosfdefault]{raleway}
\usefonttheme[onlymath]{serif}
\usepackage{microtype}

% ---------- Core packages ----------------------------------------------------
\usepackage{amsmath,amssymb}
\usepackage{booktabs}
\usepackage{tabularx,array}
\usepackage{ragged2e}
\usepackage{graphicx}
\usepackage{tikz}
\usetikzlibrary{calc,positioning,arrows.meta,shapes.geometric}
\usepackage[most]{tcolorbox}
\usepackage{etoolbox}
\usepackage{qrcode}
\usepackage{hyperref}

% ---------- Figure search paths ---------------------------------------------
\graphicspath{
  {isba2026_poster/figures/generated/}
  {isba2026_poster/figures/frozen/}
  {figures/generated/}
  {figures/frozen/}
  {Figures/multivariate_synthesis_by_cutoff/}
  {../Figures/multivariate_synthesis_by_cutoff/}
}

% ---------- Semantic palette -------------------------------------------------
\definecolor{PosterNavy}{HTML}{17324D}
\definecolor{PosterInk}{HTML}{202A33}
\definecolor{MutedGrey}{HTML}{5E6B73}
\definecolor{PanelLine}{HTML}{D8E0E3}
\definecolor{SoftPanel}{HTML}{F3F6F5}
\definecolor{SoftTeal}{HTML}{E7F2F2}
\definecolor{ExalTeal}{HTML}{006D77}
\definecolor{GloFASBlue}{HTML}{0072B2}
\definecolor{NWSOrange}{HTML}{D55E00}
\definecolor{AlGold}{HTML}{D99A00}
\definecolor{CovariateGreen}{HTML}{009E73}
\definecolor{CovariateSky}{HTML}{56B4E9}
\definecolor{GDPurple}{HTML}{7058A3}
\definecolor{WarmPaper}{HTML}{FCFDFC}

\hypersetup{
  colorlinks=true,
  urlcolor=ExalTeal,
  linkcolor=ExalTeal
}
\urlstyle{same}

% ---------- Poster metadata --------------------------------------------------
\newcommand{\PosterTitle}{Dynamic Bayesian Quantile Synthesis for Medium-Range River-Flow Forecasting}
\newcommand{\PosterAuthors}{Antonio Aguirre \quad Raquel Prado \quad Bruno Sanso}
\newcommand{\PosterAffiliation}{Department of Statistics, University of California, Santa Cruz}
\newcommand{\PosterEmail}{jaguir@ucsc.edu}
\newcommand{\PosterURL}{https://github.com/AntonioAPDL/Project1}
\newcommand{\PosterConference}{ISBA 2026 World Meeting}
\newcommand{\PosterLocation}{Nagoya, Japan}

\title{\PosterTitle}
\author{\PosterAuthors}
\institute{\PosterAffiliation}

% ---------- Page and Beamer setup -------------------------------------------
\setbeamertemplate{navigation symbols}{}
\setbeamersize{text margin left=1.85cm,text margin right=1.85cm}
\setbeamercolor{background canvas}{bg=WarmPaper}
\setbeamercolor{normal text}{fg=PosterInk,bg=WarmPaper}
\setbeamercolor{itemize item}{fg=ExalTeal}
\setbeamercolor{itemize subitem}{fg=ExalTeal}
\setbeamertemplate{itemize item}{\raisebox{0.18ex}{\tikz\fill[ExalTeal] (0,0) circle (0.11em);}}
\setbeamertemplate{itemize subitem}{\raisebox{0.18ex}{\tikz\fill[PosterNavy] (0,0) circle (0.09em);}}

\setlength{\parindent}{0pt}
\setlength{\parskip}{0.16cm}
\renewcommand{\arraystretch}{1.18}

% ---------- Type scale -------------------------------------------------------
\newcommand{\DisplayFont}{\fontfamily{Raleway-TLF}\selectfont}
\newcommand{\PosterTitleFont}{\DisplayFont\fontsize{78}{83}\selectfont}
\newcommand{\PosterAuthorFont}{\fontsize{28}{34}\selectfont}
\newcommand{\PosterMetaFont}{\fontsize{21}{26}\selectfont}
\newcommand{\PosterClaimFont}{\DisplayFont\bfseries\fontsize{32}{38}\selectfont}
\newcommand{\PosterMetricFont}{\DisplayFont\bfseries\fontsize{58}{58}\selectfont}
\newcommand{\SectionFont}{\DisplayFont\bfseries\fontsize{30}{35}\selectfont}
\newcommand{\BodyFont}{\fontsize{24}{30}\selectfont}
\newcommand{\SmallBodyFont}{\fontsize{21.5}{27}\selectfont}
\newcommand{\CaptionFont}{\fontsize{20.5}{26}\selectfont}
\newcommand{\FootFont}{\fontsize{17.5}{22}\selectfont}
\newcommand{\ChipFont}{\DisplayFont\bfseries\fontsize{19}{22}\selectfont}

% ---------- Reusable visual components --------------------------------------
\newcommand{\SectionTag}[1]{%
  \tikz[baseline=(tag.base)]{
    \node[
      fill=PosterNavy,
      text=white,
      rounded corners=1.1mm,
      inner xsep=0.34cm,
      inner ysep=0.18cm,
      font=\DisplayFont\bfseries\fontsize{17}{17}\selectfont
    ] (tag) {#1};
  }%
}

\newcommand{\SectionHead}[2]{%
  \noindent
  \begin{tabularx}{\linewidth}{@{}l@{\hspace{0.35cm}}X@{}}
    \SectionTag{#1} & {\SectionFont\color{PosterNavy}\RaggedRight #2}
  \end{tabularx}
  \vspace{0.10cm}
  {\color{PanelLine}\rule{\linewidth}{0.65pt}}\par
  \vspace{0.34cm}
}

\newtcolorbox{QuestionBox}{
  enhanced,
  colback=SoftTeal,
  colframe=SoftTeal,
  boxrule=0pt,
  arc=1.6mm,
  left=0.55cm,right=0.55cm,top=0.42cm,bottom=0.42cm,
  borderline west={2.4mm}{0pt}{ExalTeal}
}

\newtcolorbox{SoftCallout}[1][]{
  enhanced,
  colback=SoftPanel,
  colframe=SoftPanel,
  boxrule=0pt,
  arc=1.5mm,
  left=0.50cm,right=0.50cm,top=0.38cm,bottom=0.38cm,
  borderline west={1.8mm}{0pt}{ExalTeal},
  #1
}

\newtcolorbox{GoldCallout}[1][]{
  enhanced,
  colback=AlGold!9,
  colframe=AlGold!9,
  boxrule=0pt,
  arc=1.5mm,
  left=0.50cm,right=0.50cm,top=0.38cm,bottom=0.38cm,
  borderline west={1.8mm}{0pt}{AlGold},
  #1
}

\newtcolorbox{HeroResult}{
  enhanced,
  colback=white,
  colframe=ExalTeal,
  boxrule=0.85pt,
  arc=1.9mm,
  left=0.62cm,right=0.62cm,top=0.62cm,bottom=0.58cm,
  borderline north={2.7mm}{0pt}{ExalTeal}
}

\newtcolorbox{FooterBox}{
  enhanced,
  colback=SoftPanel,
  colframe=SoftPanel,
  boxrule=0pt,
  arc=1.4mm,
  left=0.55cm,right=0.55cm,top=0.38cm,bottom=0.38cm
}

\newcommand{\Pill}[3][white]{%
  \tikz[baseline=(pill.base)]{
    \node[
      fill=#2,
      text=#1,
      rounded corners=1.7mm,
      inner xsep=0.34cm,
      inner ysep=0.16cm,
      font=\ChipFont
    ] (pill) {#3};
  }%
}

\newcommand{\FactPill}[1]{%
  \tikz[baseline=(pill.base)]{
    \node[
      draw=PanelLine,
      line width=0.65pt,
      fill=white,
      text=PosterNavy,
      rounded corners=1.6mm,
      inner xsep=0.32cm,
      inner ysep=0.15cm,
      font=\DisplayFont\bfseries\fontsize{17.5}{20}\selectfont
    ] (pill) {#1};
  }%
}

\newcommand{\SourceCard}[3]{%
  \begin{tcolorbox}[
    enhanced,
    colback=white,
    colframe=PanelLine,
    boxrule=0.55pt,
    arc=1.4mm,
    height=4.85cm,
    valign=top,
    left=0.42cm,right=0.38cm,top=0.34cm,bottom=0.25cm,
    borderline west={1.7mm}{0pt}{#1}
  ]
  {\SmallBodyFont\bfseries\color{PosterNavy}#2\par}
  \vspace{0.05cm}
  {\CaptionFont\color{MutedGrey}#3\par}
  \end{tcolorbox}%
}

\newcommand{\TakeawayCard}[3]{%
  \begin{tcolorbox}[
    enhanced,
    colback=white,
    colframe=PanelLine,
    boxrule=0.55pt,
    arc=1.5mm,
    height=7.4cm,
    valign=top,
    left=0.55cm,right=0.50cm,top=0.45cm,bottom=0.35cm,
    borderline north={1.5mm}{0pt}{#1}
  ]
  \Pill{#1}{#2}\par
  \vspace{0.24cm}
  {\SmallBodyFont\color{PosterInk}#3\par}
  \end{tcolorbox}%
}

\newenvironment{PosterItems}{%
  \begin{itemize}
  \setlength{\itemsep}{0.28cm}
  \setlength{\parskip}{0pt}
  \setlength{\parsep}{0pt}
}{\end{itemize}}

\newcommand{\FigureCaption}[2][]{%
  \vspace{0.20cm}
  {\CaptionFont\color{MutedGrey}\RaggedRight
    \ifstrempty{#1}{}{\textbf{\color{PosterNavy}#1}\enspace}
    #2\par
  }
}

% ---------- Robust figure loader --------------------------------------------
% This searches the known project directories. If a figure is unavailable, the
% document still compiles with a clean placeholder that reports the filename.
\newcommand{\ResolveGraphic}[1]{%
  \gdef\ResolvedGraphic{}%
  \IfFileExists{#1}{\gdef\ResolvedGraphic{#1}}{%
  \IfFileExists{isba2026_poster/figures/generated/#1}{\gdef\ResolvedGraphic{isba2026_poster/figures/generated/#1}}{%
  \IfFileExists{isba2026_poster/figures/frozen/#1}{\gdef\ResolvedGraphic{isba2026_poster/figures/frozen/#1}}{%
  \IfFileExists{figures/generated/#1}{\gdef\ResolvedGraphic{figures/generated/#1}}{%
  \IfFileExists{figures/frozen/#1}{\gdef\ResolvedGraphic{figures/frozen/#1}}{%
  \IfFileExists{Figures/multivariate_synthesis_by_cutoff/#1}{\gdef\ResolvedGraphic{Figures/multivariate_synthesis_by_cutoff/#1}}{%
  \IfFileExists{../Figures/multivariate_synthesis_by_cutoff/#1}{\gdef\ResolvedGraphic{../Figures/multivariate_synthesis_by_cutoff/#1}}{}%
  }}}}}}%
}

\newcommand{\PosterGraphic}[3][10cm]{%
  \ResolveGraphic{#2}%
  \ifdefempty{\ResolvedGraphic}{%
    \begin{tcolorbox}[
      enhanced,
      width=\linewidth,
      height=#1,
      valign=center,
      colback=SoftPanel,
      colframe=PanelLine,
      boxrule=0.6pt,
      arc=1.3mm,
      halign=center
    ]
      {\DisplayFont\bfseries\fontsize{21}{25}\selectfont\color{PosterNavy}#3\par}
      \vspace{0.18cm}
      {\fontsize{15.5}{19}\selectfont\color{MutedGrey}\nolinkurl{#2}\par}
    \end{tcolorbox}%
  }{%
    \centering
    \includegraphics[width=\linewidth,height=#1,keepaspectratio]{\ResolvedGraphic}\par
  }%
}

% ---------- Header -----------------------------------------------------------
\newcommand{\PosterHeader}{%
  \begin{tikzpicture}[remember picture,overlay]
    \fill[SoftPanel] (current page.north west) rectangle ([yshift=-14.8cm]current page.north east);
    \fill[ExalTeal] (current page.north west) rectangle ([yshift=-0.42cm]current page.north east);
    \draw[PanelLine,line width=0.75pt]
      ([yshift=-14.8cm]current page.north west) -- ([yshift=-14.8cm]current page.north east);
  \end{tikzpicture}

  \vspace*{0.20cm}
  \begin{minipage}[t]{0.735\textwidth}
    {\PosterTitleFont\bfseries\color{PosterNavy}\RaggedRight
      Dynamic Bayesian Quantile Synthesis for\\[-0.04em]
      Medium-Range River-Flow Forecasting\par
    }
    \vspace{0.45cm}
    {\PosterAuthorFont\color{PosterInk}\PosterAuthors\par}
    \vspace{0.14cm}
    {\PosterMetaFont\color{MutedGrey}
      \PosterAffiliation\quad\textbullet\quad
      \href{mailto:\PosterEmail}{\PosterEmail}\par
    }
  \end{minipage}
  \hfill
  \begin{minipage}[t]{0.225\textwidth}
    \begin{tcolorbox}[
      enhanced,
      colback=white,
      colframe=PanelLine,
      boxrule=0.65pt,
      arc=1.6mm,
      left=0.55cm,right=0.45cm,top=0.48cm,bottom=0.45cm,
      borderline north={1.8mm}{0pt}{PosterNavy}
    ]
      {\DisplayFont\bfseries\fontsize{27}{31}\selectfont\color{PosterNavy}
        \PosterConference\par}
      \vspace{0.10cm}
      {\PosterMetaFont\color{MutedGrey}\PosterLocation\par}
      \vspace{0.35cm}
      {\color{PanelLine}\rule{\linewidth}{0.6pt}\par}
      \vspace{0.30cm}
      {\SmallBodyFont\bfseries\color{PosterInk}San Lorenzo River\par}
      {\PosterMetaFont\color{MutedGrey}Big Trees USGS gauge\par}
    \end{tcolorbox}
  \end{minipage}

  \vspace{0.50cm}
  \begin{tcolorbox}[
    enhanced,
    colback=white,
    colframe=white,
    boxrule=0pt,
    arc=1.8mm,
    left=0.60cm,right=0.65cm,top=0.44cm,bottom=0.42cm,
    borderline west={3.0mm}{0pt}{ExalTeal}
  ]
    \begin{minipage}[c]{0.115\linewidth}
      {\PosterMetricFont\color{ExalTeal}4/5\par}
      {\PosterMetaFont\bfseries\color{PosterNavy}rolling origins\par}
    \end{minipage}
    \hfill
    \begin{minipage}[c]{0.84\linewidth}
      {\PosterClaimFont\color{PosterNavy}\RaggedRight
        exAL-M-T1 has the lowest 28-day CRPS at four cutoffs;
        AL-M-T1 leads at the 2022-12-25 cutoff.\par
      }
    \end{minipage}
  \end{tcolorbox}
}

% ============================================================================
% DOCUMENT
% ============================================================================
\begin{document}
\begin{frame}[t]
\RaggedRight
\PosterHeader

\vspace{0.82cm}

{\BodyFont
\begin{columns}[T,totalwidth=\textwidth]

% ======================== LEFT COLUMN ========================================
\begin{column}{0.255\textwidth}
\begin{minipage}[t][76.2cm][t]{\linewidth}
\vspace{0pt}

\SectionHead{01}{Forecasting challenge}

\begin{QuestionBox}
{\SmallBodyFont\bfseries\color{PosterNavy}
How can heterogeneous hydrologic forecast products be reconciled into calibrated
predictive quantiles as conditions and lead times evolve?\par}
\end{QuestionBox}

\vspace{0.30cm}
{\SmallBodyFont
Medium-range decisions require a predictive distribution, but the available
information is neither exchangeable nor horizon-aligned. Forecast quality changes
with source, lead time, and hydrologic regime.\par}

\vspace{0.28cm}
\begin{columns}[T,totalwidth=\linewidth]
\begin{column}{0.49\linewidth}
  \SourceCard{PosterInk}{USGS target}{Observed daily flow for state updating and post-cutoff verification.}
\end{column}
\hfill
\begin{column}{0.49\linewidth}
  \SourceCard{GloFASBlue}{ECMWF / GloFAS}{Retrospective and operational guidance for the 28-day window.}
\end{column}
\end{columns}
\vspace{0.12cm}
\begin{columns}[T,totalwidth=\linewidth]
\begin{column}{0.49\linewidth}
  \SourceCard{NWSOrange}{NOAA / NWS}{Operational guidance on the common eight-day horizon.}
\end{column}
\hfill
\begin{column}{0.49\linewidth}
  \SourceCard{CovariateGreen}{Transfer information}{Precipitation, soil moisture, and the GDPC climate index.}
\end{column}
\end{columns}

\vspace{0.24cm}
\begin{SoftCallout}
{\SmallBodyFont
\textbf{Source-aware correction.} External products enter through
source-specific discrepancy states rather than as one homogeneous ensemble.\par}
\end{SoftCallout}

\vfill

\SectionHead{02}{Held-out rolling-origin validation}
\PosterGraphic[12.8cm]{rolling_origin_timeline.pdf}{Rolling-origin validation timeline}
\FigureCaption[Strict temporal holdout.]{At each cutoff, the model is fit using
only information available through the origin. The following 28 days are verified
against post-cutoff USGS observations; days 1--8 define the NWS-compatible comparison.}

\vspace{0.28cm}
\centering
\FactPill{5 origins}\hspace{0.16cm}
\FactPill{28-day holdout}\hspace{0.16cm}
\FactPill{NWS: days 1--8}\par
\vspace{0.16cm}
\FactPill{score scale: $\log(1+Q)$}\hspace{0.16cm}
\FactPill{lower CRPS is better}\par
\RaggedRight

\vfill

\SectionHead{03}{Model code and scoring rule}

\begin{columns}[T,totalwidth=\linewidth]
\begin{column}{0.31\linewidth}
  \centering\Pill{ExalTeal}{exAL}\par
  \vspace{0.16cm}
  {\CaptionFont\color{MutedGrey}extended likelihood family\par}
\end{column}
\begin{column}{0.31\linewidth}
  \centering\Pill{GloFASBlue}{M}\par
  \vspace{0.16cm}
  {\CaptionFont\color{MutedGrey}multivariate source synthesis\par}
\end{column}
\begin{column}{0.31\linewidth}
  \centering\Pill{AlGold}{T1}\par
  \vspace{0.16cm}
  {\CaptionFont\color{MutedGrey}retained forecast-window transfer\par}
\end{column}
\end{columns}

\vspace{0.34cm}
\begin{SoftCallout}
\centering
{\SmallBodyFont\color{PosterNavy}
\[
\operatorname{CRPS}(F,y)
=\int_{-\infty}^{\infty}\!\left(F(z)-\mathbf 1\{y\le z\}\right)^2\,dz
\]
}\vspace{-0.20cm}
{\CaptionFont\color{MutedGrey}
A proper score for the full predictive distribution; lower values indicate
better distributional forecasts.\par}
\end{SoftCallout}

\end{minipage}
\end{column}

\hfill

% ======================== CENTER COLUMN ======================================
\begin{column}{0.420\textwidth}
\begin{minipage}[t][76.2cm][t]{\linewidth}
\vspace{0pt}

\SectionHead{04}{Dynamic Bayesian quantile synthesis}
\PosterGraphic[17.2cm]{model_schematic.pdf}{Dynamic quantile synthesis schematic}
\FigureCaption[One evolving predictive system.]{For each target quantile, the
model combines multivariate forecast information with source-specific discrepancy
states, trend and seasonal dynamics, and a retained forecast-window transfer component.}

\vspace{0.26cm}
\begin{SoftCallout}
{\SmallBodyFont\color{PosterInk}
\textbf{Poster-level model summary:}
multivariate source synthesis $+$ dynamic discrepancy correction $+$
trend/seasonal states $+$ precipitation, soil moisture, and GDPC transfer
$\longrightarrow$ posterior predictive quantiles.\par}
\end{SoftCallout}

\vfill

\SectionHead{05}{Five-cutoff evidence}
\begin{HeroResult}
  \begin{minipage}[c]{0.77\linewidth}
    {\SmallBodyFont\bfseries\color{PosterNavy}
      28-day CRPS across all rolling origins\par}
  \end{minipage}
  \hfill
  \begin{minipage}[c]{0.20\linewidth}
    \RaggedLeft{\FactPill{CRPS $\downarrow$}\par}
  \end{minipage}

  \vspace{0.26cm}
  \PosterGraphic[27.6cm]{crps_28d_poster.pdf}{Five-cutoff 28-day CRPS comparison}

  \vspace{0.34cm}
  \begin{tcolorbox}[
    enhanced,
    colback=SoftTeal,
    colframe=SoftTeal,
    boxrule=0pt,
    arc=1.5mm,
    left=0.50cm,right=0.50cm,top=0.38cm,bottom=0.38cm
  ]
    \begin{minipage}[c]{0.16\linewidth}
      {\PosterMetricFont\color{ExalTeal}4/5\par}
    \end{minipage}
    \hfill
    \begin{minipage}[c]{0.80\linewidth}
      {\SmallBodyFont\color{PosterInk}
        \textbf{exAL-M-T1} has the lowest 28-day CRPS at four rolling origins
        and is retained as the extended-likelihood multivariate reference.\par}
    \end{minipage}
  \end{tcolorbox}

  \vspace{0.24cm}
  \begin{GoldCallout}
  {\SmallBodyFont
    \textbf{Important exception.} At the 2022-12-25 origin,
    \textbf{AL-M-T1} has the lowest 28-day CRPS. The evidence supports broad,
    not universal, dominance of exAL-M-T1.\par}
  \end{GoldCallout}
\end{HeroResult}

\vspace{0.28cm}
{\CaptionFont\color{MutedGrey}
Raw GloFAS provides the 28-day product reference. Raw NWS is evaluated only in
the separate horizon-aligned eight-day comparison.\par}

\end{minipage}
\end{column}

\hfill

% ======================== RIGHT COLUMN =======================================
\begin{column}{0.255\textwidth}
\begin{minipage}[t][76.2cm][t]{\linewidth}
\vspace{0pt}

\SectionHead{06}{NWS horizon nuance}
\PosterGraphic[13.2cm]{crps_8d_nws_poster.pdf}{Eight-day NWS-horizon CRPS comparison}
\FigureCaption[Horizon-aligned comparison.]{NWS contributes eight valid daily
leads at these origins. This panel is the direct NWS comparison and should not be
read as the 28-day medium-range benchmark.}

\vspace{0.28cm}
\begin{GoldCallout}
{\SmallBodyFont
\textbf{Design principle:} compare products only where their lead-time contracts
overlap. This avoids penalizing a shorter-horizon source for forecasts it does not issue.\par}
\end{GoldCallout}

\vfill

\SectionHead{07}{One forecast window, for intuition}
\PosterGraphic[18.0cm]{cutoff_2022_12_25_multivariate_synthesis_with_reference_ensembles.png}{Representative 2022-12-25 synthesis}
\FigureCaption[Illustrative, not aggregate evidence.]{Representative exAL-M-T1
synthesis at the 2022-12-25 origin with reference forecast products. The principal
validation remains the five-cutoff score comparison.}

\vspace{0.24cm}
\begin{SoftCallout}
{\SmallBodyFont
The example shows how source information and posterior uncertainty evolve across
a forecast window; it does not establish the model ranking by itself.\par}
\end{SoftCallout}

\vfill

\SectionHead{08}{Interpretation and limits}
\begin{PosterItems}
  \item The multivariate active-transfer specifications are strongest across the
  corrected rolling-origin evaluation.
  \item exAL-M-T1 is the selected extended-likelihood reference because of its
  broad five-origin performance, not because it wins every cutoff.
  \item The evaluation covers five flood-relevant origins rather than a dense,
  continuous hindcast; generalization across regimes remains a key next test.
\end{PosterItems}

\vspace{0.25cm}
\begin{QuestionBox}
{\SmallBodyFont\color{PosterInk}
\textbf{Bottom line.} Dynamic Bayesian quantile synthesis offers a coherent way
to combine heterogeneous forecast sources while preserving source discrepancy,
lead-time structure, and distributional uncertainty.\par}
\end{QuestionBox}

\end{minipage}
\end{column}

\end{columns}
}

% ======================== BOTTOM TAKEAWAYS ===================================
\vspace{0.62cm}
{\BodyFont
\SectionHead{09}{Three messages to remember}
\begin{columns}[T,totalwidth=\textwidth]
\begin{column}{0.32\textwidth}
  \TakeawayCard{ExalTeal}{1}{\textbf{Source-aware dynamics.}
  Observations, retrospective products, operational forecasts, and covariates
  enter as distinct information channels rather than one exchangeable ensemble.}
\end{column}
\hfill
\begin{column}{0.32\textwidth}
  \TakeawayCard{PosterNavy}{2}{\textbf{Validation first.}
  The main evidence is post-cutoff scoring at five rolling origins. A single
  hydrograph is explanatory context, not the validation claim.}
\end{column}
\hfill
\begin{column}{0.32\textwidth}
  \TakeawayCard{NWSOrange}{3}{\textbf{Horizon contracts matter.}
  The medium-range result uses the 28-day GloFAS-compatible window; raw NWS is
  compared only on common leads 1--8.}
\end{column}
\end{columns}
}

% ======================== FOOTER =============================================
\vspace{0.38cm}
\begin{FooterBox}
\begin{minipage}[c]{0.77\linewidth}
  {\FootFont\color{MutedGrey}\RaggedRight
  \textbf{Provenance.} Poster figures and numerical claims are generated from
  frozen manuscript-facing artifacts dated 2026-06-19. Core references:
  Yu and Moyeed (2001), Prado and West (2010), Yan et al. (2025), and the revised
  article artifact manifests.\par
  \vspace{0.16cm}
  \textbf{Contact.} \href{mailto:\PosterEmail}{\PosterEmail}\quad\textbullet\quad
  \href{\PosterURL}{github.com/AntonioAPDL/Project1}\par
  }
\end{minipage}
\hfill
\begin{minipage}[c]{0.20\linewidth}
  \begin{columns}[c,totalwidth=\linewidth]
  \begin{column}{0.42\linewidth}
    \centering
    \qrcode[height=2.75cm]{\PosterURL}
  \end{column}
  \begin{column}{0.55\linewidth}
    {\FootFont\bfseries\color{PosterNavy}Code, figures, and artifacts\par}
    \vspace{0.08cm}
    {\FootFont\color{MutedGrey}Scan for the project repository and reproducibility materials.\par}
  \end{column}
  \end{columns}
\end{minipage}
\end{FooterBox}

\end{frame}
\end{document}
